\pagestyle{empty}
\label{cover:UseCase}

\noindent \textbf{\Large CD Használati útmutató}

\vskip 1cm

A CD mellékletben található 2 mappa: a Notes és a Codes.\\
A Notes mappán belül található meg a LaTeX forráskód. A Codes mappán belül pedig a Godot projekt minden hozzátartozó elemével együtt.

\vskip 1cm

\noindent \textbf{\Large A Godot projekt beüzemelése}

\vskip 1cm

Ahhoz, hogy a Godot projektet be tudjuk üzemelni először is le kell tölteni a Godot Engine - .NET változatát. Ez a Godot hivatalos weboldalán mind elérhető \cite{Godot}. Amikor elindítjuk a játékmotort a bal felső sarokban lesz 3 opció, ezek közül az "Import" opciót kell kiválasztanunk. Importálás előtt készítsünk egy biztonsági másolatot a projektről. Ezek után navigáljunk el a \texttt{Thesis/Codes/yagster} mappához, majd válasszuk azt ki és kattintsunk a \texttt{megnyitás} vagy \texttt{open} gombra. A Godot fel fogja ismerni, mint egy projekt és az importálás sikeres lesz. Ha egy olyan üzenetet do fel, hogy a projekt egy régebbi verzióban készült nem kell megijedni, a projekt attól még importálható és futtatható. Nagy valószínűséggel nem lesz baj belőle, de erre van a biztonsági mentés. Ezek után kétszer az importált projektre kattintva a Godot meg fogja nyitni. Innen a projekt elindítható a jobb felső sarokban lévő nyíl gombbal.